\begin{scanpage}{0236}{218}
\noindent first class. As an example of such a problem, one may point to the problem of diffraction of a plane wave by a frustum of a cone or by an infinitely long cylinder with a polygonal transverse cross section. Those problems whose solution requires obtaining (using the methods of mathematical diffraction theory) a whole series of new results must be referred to the second class. In particular, in order to give a complete solution to the diffraction problem of a finite cone, it is necessary to have more precise knowledge on the diffraction laws of a semi-infinite cone.

Summing up, one may say that physical diffraction theory aids one in analyzing and sorting out the diffraction phenomena for complex bodies, poses problems for treatment by mathematical diffraction theory, and enables one to effectively apply the rigorous results of mathematical diffraction theory for the solution of new problems.

In conclusion, I express my deep thanks to L. A. Vaynshteyn for his valuable advice and regular discussion of the questions to which this book is devoted, and also for his attentive reading of the manuscript and for a number of useful remarks. I also take this opportunity to express sincere thanks to M. L. Levin for his interest in this work and his helpful remarks.
\end{scanpage}
